\paragraph{Literal-time multiplicity bounds on the seven--five family.}
For $N=7\cdot5^e$, $e\ge1$, every eigenvalue of the actual $e^{itA_N}M_N$
has multiplicity at most seven at every nonzero real algebraic time,
including $t=10^{-80}N$. At least $612N/625+7$ eigenvalues are simple for even
$e\ge4$, and at least $612N/625+175$ for even $e\ge6$; the odd counts are $24N/25$ for $e\ge3$,
$24N/25+35$ for $e\ge5$, and $24N/25+875$ for $e\ge7$. The exact special-dimension lower bounds are
4333 at $N=4375$ and 107443 at $N=109375$.
The separately reviewed cat, complete oldform, deepest-core and newest-sector
arguments supply the local splitting premises. Polynomial minors in $U(t)$
are algebraic exponential polynomials, so the registered Lindemann--Weierstrass
input transfers witnessed ranks to each nonzero algebraic real time.
The rank of $\chi'(U)$ counts simple eigenvalues; higher derivative stacks
recover the dimensions of multiplicity-bounded eigenspaces.

A finite first-order companion proves simple $U$ at $N=35,175$; the latter
is the additional bounded case. Zero oldform first-order compressions defeat
a criterion, not simplicity of the actual propagator. The repeated-eigenspace
count does not assert persistence of a fixed oldform invariant subspace.
Full uniform simplicity, quantitative gaps and BGS remain open.

The underlying inert-torus multiplicity-free mechanism has prior provenance
in Kelmer~\cite{Kelmer2008PrimePowers}; the actual cyclic/CRT and jet arguments
are independently derived here.

% Source: mathematics/literal_time_multiplicity_bounds.md and its three
% source-bound intermediate-old-layer components. FM-026 only.
\paragraph{Intermediate old layers and the strengthened partial count.}
For the unchanged family $N=7\cdot5^e$ and every nonzero real algebraic
$t$, the preceding multiplicity bound remains at most seven. The exact
first-old-layer classifier strengthens the simple-eigenvalue lower bound to
$612N/625+7$ for even $e\ge4$. Its regular norm classes contribute half of
the first-old quotient; the singular classes are decided by an exact
cyclotomic coefficient-balancing test. Nonempty intersection alone does
not ensure a nonzero compression: at residual dimension $625$, the norm
label $c=13$ has five intersections but six zero blocks and 24 simple blocks.
The general return operator is a character-normalized matrix inverse with
initial factors and the resonant projector term retained. Its corrected
coherent-lattice residue is an affine orbit sum, and integral circlewise
six-cycle bases prove second-order splitting of all 24 intermediate
seven-dimensional blocks at residual dimension $25$. Including the prior
$3,2,2$ deepest blocks gives $612N/625+175$ simple eigenvalues for even
$e\ge6$. Exact first-old counts give 4333 at $N=4375$ and 107443 at
$N=109375$. The counts add over disjoint complete unperturbed cat eigenspaces
and pass to algebraic time through the existing exponential-rank theorem
and its registered Lindemann--Weierstrass input. They do not identify
fixed-time eigenvectors with old invariant subspaces. Larger intermediate
residuals and their wild-character denominators remain open; no full
simplicity, gap estimate, or BGS conclusion follows.

The separately audited odd newest sector contributes $24N/25$ first-order
branches at odd depths $e\ge3$. The deepest odd core adds $35$ for $e\ge5$.
At odd $e\ge7$, the residual-$125$ bridge supplies $840$ further second-order
branches; nesting places these and the deepest core in $W_1$, disjoint from
the newest complete cat groups. Thus the addend is $875=840+35$, counted once.
The inherited algebraic-time rank transfer applies; no fixed-time oldform
invariance, full simplicity or quantitative gap is asserted.

